\chapter{Reinforcement learning and curriculum (concept)}

\section{MDP framing}
\begin{itemize}
  \item \textbf{Environment:} Showdown-backed battle through poke-env and the Gymnasium wrapper.
  \item \textbf{State:} Encoded battle observations (token grid, categorical IDs, action mask).
  \item \textbf{Actions:} Discrete move/switch indices (including format-specific gimmick-expanded slots where applicable).
  \item \textbf{Reward:} Sum of terminal outcome terms, HP/faint shaping, optional step penalties and matchup-related shaping (see \texttt{RewardConfig}).
\end{itemize}

\section{Algorithm}
Training uses \textbf{Proximal Policy Optimization (PPO)} via Ray RLlib with a custom PyTorch policy/value network (transformer trunk).

\section{Curriculum (conceptual)}
Because the full challenge is harsh, training progresses through stages with different opponent mixtures and sometimes different reward weights. \textbf{Stage promotion} is driven by a rolling win-rate window over completed episodes (\texttt{CurriculumManager} in \texttt{src/training/curriculum.py}). When a stage changes, the environment bridge applies a payload that updates both \textbf{opponent sampling} and \textbf{reward configuration}---not a cosmetic toggle.

\section{Staged scenarios (original planning ideas)}
Early project notes also listed hypothetical curricula (single-Pok\'{e}mon battles, small teams, full Nuzlocke rules, full challenge graph). The \textbf{implemented} default stages are enumerated in the next chapter.
